\chapter{Saint-Gobain Sekurit et le contexte industriel}

\section{Présentation du site de Kénitra}

Saint-Gobain Sekurit fabrique des vitrages automobiles pour la première monte
et le marché du remplacement. Implanté dans l'Atlantic Free Zone, le site de
Kénitra produit des pare-brise automobiles et dessert plusieurs constructeurs.
Son organisation associe les opérations de transformation du verre, le
marquage, l'assemblage et le contrôle qualité.

\begin{figure}[H]
  \centering
  \includegraphics[width=0.68\textwidth]{usine-kenitra.jpg}
  \caption{Vue du site Saint-Gobain Sekurit de Kénitra}
\end{figure}

\begin{table}[H]
\centering
\caption{Fiche synthétique de l'organisme d'accueil}
\begin{tabularx}{0.93\textwidth}{>{\bfseries}p{0.31\textwidth}Y}
\toprule
Élément & Présentation synthétique \\
\midrule
Entreprise & Saint-Gobain Sekurit Maroc \\
Implantation & Atlantic Free Zone, Kénitra \\
Activité & Production de pare-brise automobiles \\
Clients cités & Stellantis, Renault, Volkswagen et Seat \\
Certifications citées & IATF 16949, ISO 9001, ISO 14001 et ISO 45001 \\
Produits présentés & Pare-brise feuilleté, PVB chauffant et PVB antenne \\
\bottomrule
\end{tabularx}
\end{table}

\section{Le pare-brise feuilleté}

Un pare-brise feuilleté associe deux feuilles de verre et un film intermédiaire
en PVB. Ce film maintient les fragments en cas de choc et peut recevoir des
fonctions supplémentaires, comme le chauffage ou une antenne. La fabrication
ne consiste donc pas seulement à découper une forme. Elle enchaîne des
opérations de préparation, d'impression, de bombage, d'assemblage et de
contrôle~\cite{sekuritprocess}.

\begin{figure}[H]
  \centering
  \begin{subfigure}[b]{0.44\textwidth}
    \centering
    \includegraphics[width=\linewidth,height=5.2cm,keepaspectratio]{pare-brise-feuillete.jpg}
    \caption{Pare-brise feuilleté}
  \end{subfigure}
  \hfill
  \begin{subfigure}[b]{0.50\textwidth}
    \centering
    \includegraphics[width=\linewidth,height=5.2cm,keepaspectratio]{procede-verre-feuillete.jpg}
    \caption{Vue d'ensemble du procédé}
  \end{subfigure}
  \caption{Produit et chaîne de fabrication du vitrage automobile}
\end{figure}

\section{Place de l'impression et du contrôle}

Dans la partie pré-processing, le verre est chargé, découpé, façonné, lavé et
imprimé. La sérigraphie du logo et de la bande noire est suivie d'un contrôle
visuel et d'un contrôle automatique. Le projet GlassInk Monitor concerne un
autre équipement d'impression, une
imprimante CIJ dédiée au marquage DMC, mais les enjeux restent proches :
position correcte, contraste suffisant, lisibilité et traçabilité.

\begin{figure}[H]
  \centering
  \begin{subfigure}[b]{0.48\textwidth}
    \centering
    \includegraphics[width=\textwidth]{serigraphie.jpg}
    \caption{Table d'impression sur verre}
  \end{subfigure}
  \hfill
  \begin{subfigure}[b]{0.48\textwidth}
    \centering
    \includegraphics[width=\textwidth]{pare-brise-controle.jpg}
    \caption{Pare-brise au poste de contrôle}
  \end{subfigure}
  \caption{Exemples d'impression et de contrôle dans la ligne}
\end{figure}

Un défaut de marquage peut avoir plusieurs origines. La donnée imprimée peut
être fausse, la machine peut manquer de fluide, la viscosité peut dériver ou la
buse peut être partiellement obstruée. Une supervision utile doit donc réunir
les informations du procédé, de l'imprimante et des consommables.

\section{Systèmes industriels présents}

Les PLC coordonnent les équipements de la ligne et échangent leurs données sur
PROFINET. Litmus Edge récupère les tags industriels, les normalise et peut les
publier vers des systèmes informatiques. Il ne remplace pas les automates :
les automatismes restent responsables du contrôle de la production.

Ce découpage fixe la frontière du projet. GlassInk Monitor commence au niveau
de la sortie MQTT de Litmus Edge. La configuration des drivers PLC et la
logique de commande de la ligne ne font pas partie du travail.
